% SPDX-License-Identifier: CC-BY-4.0
% Copyright (C) 2026 Fabian Januszewski, Abhijit (the Langlandscape team)
% Stage 06 section of the submission report (proof outlines condensed from mathematics/chapters/s06; statements verbatim).
% In the verbatim boxes, the chapter's \lit{stem}{loc} literature pointers were mechanically converted to \cite[loc]{KEY}
% (no other change to the quoted text). No quoted box contained a person or node name; Remark 23.8 is paraphrased, not quoted.
\subsection{Stage 06: the integral Jacquet--Langlands defect}\label{ss:stage06}

Stage 06 is the end point of the route. On the $\GL_2$ side, with the definite algebra $B=B_q$ (ramified at $\{q,\infty\}$) entering
only through the $\Lam$-adic Hecke Jacquet--Langlands isomorphism $\T'\cong\T^{\ord}_B$ of stages 01 and 04, it proves that the
defect $(h_F)=\eta_F(\eta^B_F)^{-1}$ of a $q$-new Hida branch $F$ equals the old/new congruence ideal $\cc_q(F)$ (pure commutative
algebra), and that $\cc_q(F)$ equals the Tamagawa ideal $\tq_q(F)$ (Galois theory plus one fixed-weight input). Specialised to an
arithmetic point this is the integral identity
\[\eta_f(N)=\cc_q(f)\cdot\eta_f(M,q),\qquad \cc_q(f)=(\varpi^{t_f(q)}),\]
where $\eta_f(M,q)$ is the congruence ideal of the $q$-new quotient, i.e.\ of the definite side via JL. The stage occupies
\S\S19--25 of main.pdf, pp.~159--188: commutative algebra (\S19), the setting and the old/new decomposition (\S20), the defect formula
(\S21), the Galois half (\S22), the case $q\equiv-1\bmod p$ (\S23), the Betti incarnation of $\cc_q$ (\S24), open cases (\S25).

\subsubsection*{Key definitions}

\emph{Gorenstein algebras and congruence ideals} (main.pdf, Def.~19.1, p.~159). $A$ is a commutative ring, $S$ a commutative
$A$-algebra, $S^\vee=\Hom_A(S,A)$ with $(s\varphi)(x)=\varphi(sx)$. An $\ell\in S^\vee$ is a \emph{trace} if $S$ is finite free over $A$
and $s\mapsto\ell(s\,\cdot)$, $S\to S^\vee$, is bijective; $S$ is \emph{$A$-Gorenstein} if it has a trace. For an $A$-algebra map
$\mu:S\to A$, $S[\mu]:=\Ann_S(\ker\mu)$ (the $\mu$-eigenspace of $S$ acting on itself) and the \emph{congruence ideal} is
$\eta(S,\mu):=\mu(S[\mu])\subseteq A$. (Lean: \texttt{IsGorensteinWith}, \texttt{IsGorenstein}, \texttt{congruenceIdeal}.)

\emph{The Hecke rings and the defect} (main.pdf, Def.~20.3, p.~166). $\Lam=\cOE[[1+p\Zp]]$, $N=qM$. $\T$ is Hida's ordinary
cuspidal Hecke algebra of tame level $\Gamma_0(N)$ (with all $U_\ell$, $\ell\mid Np$), localised at a non-Eisenstein $\m$; $\T'$, $\T''$
are its images on the $q$-new and the $q$-old ordinary $\Lam$-adic cusp forms; $\fa=\ker(\T\to\T')$, $\fb=\ker(\T\to\T'')$. For the
primitive $q$-new branch $\lambda_F:\T\to I_F$, put $R=\T\otimes_\Lam I_F$, $R'=\T'\otimes_\Lam I_F$, $R''=\T''\otimes_\Lam I_F$,
$\lambda=\lambda'\circ(R\to R')$, and
\[\eta_F:=\eta(R,\lambda),\qquad\eta^B_F:=\eta(R',\lambda'),\qquad(h_F):=\eta_F(\eta^B_F)^{-1}\ \text{ when both are principal.}\]
The seed's $h_F=\eta_F/\eta^B_F$ makes sense only under (Gor$_R$), (Gor$_{R'}$) and $\eta^B_F\ne0$.

\emph{Old/new congruence ideal} (main.pdf, Def.~20.5, p.~168). $\cc:=\mathrm{im}(\fb\to R')=(\fa+\fb)/\fa$, so $R'/\cc\cong R/(\fa+\fb)$
is Diamond's $\Lam$-adic old/new congruence module \cite{Dia89}; $\cc_q(F):=\lambda'(\cc)=\lambda(\fb)\subseteq I_F$. At an arithmetic
point with $q$-new specialisation $f$: $\cc_q(f):=\lambda_f(\fb_k)$, $\fb_k=\ker(\T_k(N)_\m\to\T''_k)$.

\emph{Tamagawa ideal} (main.pdf, Def.~20.6, p.~169). $H$ is Hida's ordinary $\Lam$-adic étale cohomology of tame level $N$ at $\m$
(free of rank two over $\T$ under (H2)); $\mathcal T_F=H\otimes_{\T,\lambda_F}I_F$ gives $\rho_F:G_\Q\to\GL_2(I_F)$, and $\tq_q(F)$ is the
ideal generated by the entries of $\rho_F(\sigma)-1$, $\sigma\in I_q$. At a point, $t_f(q)$ is the Tamagawa exponent of
\cite[Def.~2.12]{KO} ($=$\cite[Def.~3.3]{PW}): the largest $t$ with $T_f/\varpi^tT_f$ unramified at $q$.

\emph{The level hypothesis} (Lev$'_M$) (main.pdf, Hyp.~23.1, p.~175), used only in \S23. With $\Sigma^-=\{\ell\mid M:\bar\rho$ unramified
at $\ell\}$, $M_0=\prod_{\ell\in\Sigma^-}\ell$ and $P_{\bar\rho}$ the set of \cite[(2.1)]{KO}: (L1) $\ell\,\|\,M$ for $\ell\in\Sigma^-$; (L2)
$\rho_f$ is \cite[Def.~3.5]{KO}-minimally ramified at every $\ell\mid N(\bar\rho)$, $\ell\notin P_{\bar\rho}$; (L3) $\bar\rho$ is
twist-minimal. (L2), (L3) are empty when $N(\bar\rho)=1$, and then $M=M_0$ is squarefree.

\emph{The standing hypotheses.} Verbatim:

\begin{quoted}{Hypothesis}{20.1}{165}{hyp:s06-H}{standing hypotheses; seed \S9.1, corrected}{chapters/s06/s06-m2.tex}\label{rpt:hyp:s06-H}
\leavevmode
\begin{itemize}
\item[(H0)] $p\ge5$, $q\ne p$, $(M,pq)=1$, $N=qM$, $p\nmid\varphi(N)$. In particular $q\not\equiv1\bmod p$.
\item[(H0$^-$)] (\emph{new}; needed for the fixed-weight input Cor.~21.2 and for Prop.~24.1;
  Thm~21.5 needs (H0$^-$) \emph{or} (Lev); the inclusion of Prop.~22.1 uses neither)
  $q\not\equiv-1\bmod p$, or $\bar\rho_\m$ is ramified at $q$. Together with (H0): $p\nmid q^2-1$, or $\bar\rho_\m$ is ramified at $q$.
\item[(Lev)] (\emph{new}; alternative to (H0$^-$) in Thm~21.5) At the point $P$ of (P), of weight $k=k_P$, with
  $f=F_P$ (ordinary): either (F0) $M=1$; or (F1) $M>1$ is squarefree, $\bar\rho$ is unramified at every $\ell\mid M$ (so $N(\bar\rho)=1$),
  and $\ell\not\equiv-1\bmod p$ for every $\ell\mid M$; or (F1$'$) as (F1), but with some $\ell\mid M$, $\ell\equiv-1\bmod p$. These are the
  cases (F0), (F1), (F1$'$) of Prop.~23.2, in which its hypothesis (Lev$'_M$) (with (L3)) holds automatically. Status:
  Prop.~23.2 is \stpmc{} in all three cases, for (F1$'$) in our sign-fixed reading R4 of \cite[Thm~3.14]{KO} (ours, not
  \cite{KO}'s text). Not included: $N(\bar\rho)>1$ and non-squarefree $M_0$, which are \stopen{} (Remark~23.8); (L3) is the condition under
  which the method may extend.

\item[(H1)] $\T$ is Hida's full ordinary cuspidal Hecke algebra of tame level $\Gamma_0(N)$ over $\Lam$, in a fixed weight component
  $k_0\bmod(p-1)$. It is generated by all $T_\ell$ ($\ell\nmid Np$), $U_\ell$ ($\ell\mid Np$, in particular $U_q$ and $U_p$), and it
  is localised at a non-Eisenstein maximal ideal $\m$. $F$ is a primitive ordinary $q$-new branch, with $\lambda_F:\T\to I_F$, and
  $I_F$ is a normal local domain finite over $\Lam$. Then $I_F$ is $\Lam$-free: it is normal of dimension $2$, hence
  Cohen--Macaulay (Serre's criterion), hence free over the regular local ring $\Lam$ (Auslander--Buchsbaum). So the seed's
  ``hence $\Lam$-free'' is correct; it is used only through flatness of $R=\T\otimes_\Lam I_F$ over $I_F$.
\item[(H2)] $\bar\rho=\bar\rho_\m$ is absolutely irreducible on $G_{\Q(\sqrt{p^*})}$ and $p$-distinguished.
\item[(Im)] $p\mid\#\bar\rho(G_\Q)$. \emph{None of the cited theorems below uses (Im).} \cite[Thm~1.1]{KO} needs only (H2) (its
  condition (2) $=$ Assumption 2.1(TW)), and \cite[Rem.~2.10]{KO} says that the \cite{DFG} and \cite{Dim09} image conditions are interchangeable
  over $\Q$. We keep (Im) as in the seed, flagged as probably superfluous (Risk R4).
\item[(P)] Either (P1) $k_0\not\equiv2\bmod(p-1)$, and then $F$ has an arithmetic point $P$ of weight $k_P\in[3,p-1]$ with trivial
  character (a \emph{Kim--Ota point}). Or (P2) $k_0\equiv2$, $F$ has a weight-two point $P$ which is the ordinary
  $p$-stabilisation of a newform $f$ of level $N$, \emph{and} (\emph{new}) the weight-two ordinary Hecke algebra
  $\T^{\ord}_2(\Gamma_0(Np))_\m$ has no $p$-new eigenforms. The extra condition makes $\T/P\T\cong\T_2(\Gamma_0(N))_\m$,
  which is the algebra of \cite{KO}. Without it the control step (3) of the proof compares with an algebra that \cite{KO} does not treat.
\end{itemize}
\sthyp
\end{quoted}

\noindent By (main.pdf, Remark~20.2, p.~166), (H2)$+$(H0) imply Hsieh's residual condition (CR,$\{q\}$) \cite[\S1.4, p.~4]{Hs}, whose congruence condition is $q\equiv+1$ only; so the
combinatorial half needs only (H0), (H1), (H2), and $q\equiv-1$ enters only through the fixed-weight Galois input.

\subsubsection*{Main results and proof outlines}

\paragraph{A. The commutative-algebra half (main.pdf, \S19).} Pure module theory over an arbitrary base $A$; the arithmetic enters only
through the hypotheses. Every statement of \S19 is formalized in the namespace \pathtt{Langlandscape.CommAlg} (sorry-free, standard
axioms only), with the qualifiers stated in each box; no \texttt{Hyp.*} is used in this layer.

\begin{quoted}{Lemma}{19.2}{160}{lem:s06-cyclic}{cyclic dual}{chapters/s06/s06-m1.tex}\label{rpt:lem:s06-cyclic}
Let $S$ be finite free over $A$. Then $S$ is $A$-Gorenstein iff $S^\vee$ is a cyclic $S$-module, and then every generator of $S^\vee$ is a
trace. \stproved{} \lean{Langlandscape.CommAlg.isGorenstein\_iff\_cyclic}; the clause ``every generator is a trace'' is \lean{Langlandscape.CommAlg.IsGorensteinWith.of\_surjective}. \stformalized{}.
\end{quoted}
\begin{outline}
If $S^\vee=S\ell$, then $s\mapsto s\ell$ is a surjection between free $A$-modules of the same finite rank; composed with an
$A$-isomorphism $S^\vee\cong S$ it is a surjective endomorphism of a finitely generated module, hence injective (Vasconcelos; the
Orzech property, Lean \pathtt{OrzechProperty.bijective\_of\_surjective\_of\_injective}). \stproved, \stformal.
\end{outline}

\begin{quoted}{Lemma}{19.3}{160}{lem:s06-S}{base change of congruence ideals; seed Lemma S}{chapters/s06/s06-m1.tex}\label{rpt:lem:s06-S}
Let $S$ be finite free and $A$-Gorenstein with trace $\ell$, and let $\mu:S\to A$ be an $A$-algebra map. Then
$S[\mu]=\Ann_S\ker\mu=A\theta$ is free of rank one, where $\theta$ is the unique element with $\ell(\theta x)=\mu(x)$ for all $x\in S$.
Hence $\eta(S,\mu)=(\mu(\theta))$ is principal. For any ring homomorphism $\varphi:A\to A'$ (not necessarily local), $S\otimes_AA'$ is $A'$-Gorenstein and
$\eta(S\otimes A',\mu\otimes A')=\varphi(\eta(S,\mu))A'$. \stproved{} \lean{Langlandscape.CommAlg.LemmaS.ann\_ker\_eq\_span},
\lean{Langlandscape.CommAlg.LemmaS.trace\_theta\_mul} (with \pathtt{theta\_unique}, \pathtt{smul\_theta\_injective}), \lean{Langlandscape.CommAlg.LemmaS.congruenceIdeal\_eq\_span};
base change \lean{Langlandscape.CommAlg.IsGorenstein.baseChange}, \lean{Langlandscape.CommAlg.LemmaS.congruenceIdeal\_baseChange}
(for an arbitrary $A$-algebra $A'$, i.e.\ $\varphi$ not necessarily local). \stformalized{}, full statement.
\end{quoted}
\begin{outline}
(1) $\mu\in S^\vee$, so bijectivity of $s\mapsto\ell(s\,\cdot)$ gives a unique $\theta$ with $\ell(\theta\,\cdot)=\mu$.
(2) $\theta\in S[\mu]$: for $x\in\ker\mu$, $\ell(\theta xs)=\mu(x)\mu(s)=0$ for all $s$, so $\theta x=0$ by injectivity.
(3) $S[\mu]\subseteq A\theta$: for $s\in S[\mu]$, $\ell(s\,\cdot)=\ell(s)\mu=\ell(\ell(s)\theta\,\cdot)$, so $s=\ell(s)\theta$.
(4) $A\theta\cong A$: $a\theta=0$ gives $a\mu=0$, so $a=a\mu(1)=0$.
(5) Base change: for $S$ finite free, $\Hom_{A'}(S\otimes A',A')=A'\otimes S^\vee$, so $\mathrm{id}\otimes\ell$ is a trace, and $1\otimes\theta$ is
the new $\theta$. Locality of $\varphi$ is not used; what is used is the existence of $\theta$, i.e.\ (Gor$_S$): without it $\Ann$ does not
commute with non-flat base change. All steps \stproved, \stformal.
\end{outline}

\begin{quoted}{Lemma}{19.4}{160}{lem:s06-G}{Gorenstein from one fibre; seed Lemma G}{chapters/s06/s06-m1.tex}\label{rpt:lem:s06-G}
Let $A$ be local, $S$ finite free over $A$ and $\mathfrak p\subseteq\m_A$ an ideal. If $S/\mathfrak pS$ is $A/\mathfrak p$-Gorenstein, then
$S$ is $A$-Gorenstein. \stproved{} \lean{Langlandscape.CommAlg.isGorenstein\_of\_isGorenstein\_quotient}. \stformalized{}, full statement.
\end{quoted}
\begin{outline}
Lift a trace $\bar\ell$ of $S/\mathfrak pS$ to $\ell\in S^\vee$ (possible since $S^\vee\otimes A/\mathfrak p=(S/\mathfrak pS)^\vee$ for $S$ finite free).
The finite $A$-module $N=S^\vee/S\ell$ satisfies $N=\mathfrak pN\subseteq\m_AN$, so $N=0$ by Nakayama over $A$ (no locality of $S$ is needed), and
Lemma~19.2 concludes. \stproved, \stformal. The chapter does not apply Lemma~19.4: the lift of the Gorenstein property from one weight to
$\Lam$ is cited from \cite{Hs}, \cite{MW86} instead.
\end{outline}

\begin{quoted}{Lemma}{19.5}{161}{lem:s06-D}{new/old duality; seed Lemma D, corrected}{chapters/s06/s06-m1.tex}\label{rpt:lem:s06-D}
Let $A$ be a Noetherian local domain and $S$ a finite free $A$-algebra, with ideals $\fa,\fb$ such that $S/\fa$, $S/\fb$ are
$A$-torsion-free, $\fa\cap\fb=0$ and $S_K=(S/\fa)_K\times(S/\fb)_K$, where $K=\mathrm{Frac}\,A$. Then:
\begin{enumerate}[label=(\alph*)]
\item $\Ann_S(\fa)=\fb$ and $\Ann_S(\fb)=\fa$. (No Gorenstein hypothesis is needed.)
\item If $S$ is $A$-Gorenstein with trace $\ell$, then $\beta\mapsto\ell(\beta\,\cdot)$ gives $S$-isomorphisms
  $\fb\cong(S/\fa)^\vee$ and $\fa\cong(S/\fb)^\vee$.
\item If moreover $S/\fb$ is $A$-free, then $\fb$ is a direct summand of $S$ as an $A$-module, hence $A$-free, and
  $s\mapsto\ell(s\,\cdot)|_\fb$ induces $S/\fa\cong\fb^\vee$. In particular $S/\fa$ is $A$-free.
\end{enumerate}
All of this commutes with base change $A\to A'$ that preserves the hypotheses. \stproved{}
(Lean: the interface \texttt{if:s06-lemD} below, \lean{Langlandscape.CommAlg.NewOldData.ann\_a}, \lean{Langlandscape.CommAlg.NewOldData.bEquivDualQuot}, \lean{Langlandscape.CommAlg.NewOldData.quotAEquivDualB}, \dots.)
\stformalized{} for (a)--(c), with these qualifiers: the hypothesis $S_K=(S/\fa)_K\times(S/\fb)_K$ enters Lean in the field-free form \texttt{split} (Lemma~19.6; the passage to the product decomposition is not formalized); the $S$-isomorphisms of (b), (c) are formalized as $A$-linear isomorphisms plus the equivariance lemmas \texttt{\_smul}; the final sentence on base change is \emph{not} formalized.
\end{quoted}

\begin{quoted}{Lemma}{19.6}{161}{lem:s06-split}{the field-free splitting}{chapters/s06/s06-m1.tex}\label{rpt:lem:s06-split}
Let $A$ be a domain with fraction field $K$ and $\fa,\fb\subseteq S$ ideals with $\fa\cap\fb=0$. Then
$S_K\to(S/\fa)_K\times(S/\fb)_K$ is an isomorphism iff $(\fa+\fb)S_K=S_K$ iff there is $c\in A\smallsetminus\{0\}$ with
$c\cdot1\in\fa+\fb$. \stproved{} \lean{Langlandscape.CommAlg.NewOldData.split\_iff}, \stformalized{} \emph{for the second equivalence only} ($\exists c\ne0$, $c\cdot1\in\fa+\fb$ $\iff$ $(\fa+\fb)S_K=S_K$, with $S_K=\mathrm{Frac}A\otimes_AS$); the first equivalence (Chinese remainder theorem, product decomposition) is not formalized.
\end{quoted}
\begin{outline}(Lemmas 19.6 and 19.5.)
Lemma~19.6: localisation is exact, so $\fa_K\cap\fb_K=0$; by the Chinese remainder theorem $S_K\to S_K/\fa_K\times S_K/\fb_K$ is injective and is
surjective iff $\fa_K+\fb_K=S_K$; clear denominators. \stproved. For Lemma~19.5 fix $c\ne0$ and $x\in\fa$, $y\in\fb$ with $c\cdot1=x+y$.
\begin{enumerate}[label=(\arabic*)]
\item (a) $\fb\fa=0$ gives $\fb\subseteq\Ann(\fa)$. Conversely, for $s\in\Ann(\fa)$, $cs=ys\in\fb$, and torsion-freeness of $S/\fb$ gives $s\in\fb$.
  Only the \texttt{NewOldData} fields and ``$A$ a domain'' are used. \stproved
\item (b) For $\beta\in\fb$, $\Phi_\fb(\beta):\bar s\mapsto\ell(\beta s)$ is well defined on $S/\fa$ by (a). Injective because $\ell$ is a trace.
  Surjective: for $\psi\in(S/\fa)^\vee$, $\psi\circ\pi_\fa=\ell(\beta\,\cdot)$ for some $\beta\in S$; then $\beta\fa=0$, so $\beta\in\Ann(\fa)=\fb$. \stproved
\item (c) $0\to\fb\to S\to S/\fb\to0$ splits since $S/\fb$ is free; so $\fb$ is f.g.\ projective, hence free over the local ring $A$.
  $\Psi:S/\fa\to\fb^\vee$, $\Psi(\bar s)(\beta)=\ell(s\beta)$, is injective because $\Ann(\fb)=\fa$ and surjective via the projection onto $\fb$;
  so $S/\fa\cong\fb^\vee$ is free. \stproved
\end{enumerate}
The seed stated (c) without a reason for the freeness of $S/\fa$; freeness of $(S/\fa)^\vee\cong\fb$ alone does not give it when $\dim A\ge2$,
and the argument uses the summand property instead. \stformal{} with the qualifiers in the box.
\end{outline}

\noindent The combinatorial half (seed Theorem~E):

\begin{quoted}{Theorem}{19.8}{163}{thm:s06-E}{the combinatorial half; seed Thm~E, hypotheses explicit}{chapters/s06/s06-m1.tex}\label{rpt:thm:s06-E}
Let $S,\fa,\fb$ be as in Lemma~19.5. Put $S'=S/\fa$, and let $\lambda:S\to A$ factor through $\lambda':S'\to A$. Assume
\begin{itemize}
\item[(Gor$_S$)] $S$ is $A$-Gorenstein.
\end{itemize}
Let $\theta$ generate $S[\lambda]$ (Lemma~19.3), and let $x\in S'$ be its image. Then:
\begin{enumerate}[label=(\roman*)]
\item $\theta\in\fb$, $x\in S'[\lambda']$, $\eta(S,\lambda)=(\lambda'(x))$, and $\cc_q:=\lambda(\fb)=\{\psi(x):\psi\in S'^\vee\}$ (the
  \emph{order ideal} of $x$).
\item If $S'[\lambda']$ is free of rank one with a generator $\theta'$ that is unimodular in $S'$ (i.e.\ $S'[\lambda']$ is an
  $A$-direct summand of $S'$), then $\eta(S,\lambda)=\cc_q\cdot\eta(S',\lambda')$.
\item This holds in particular under (Gor$_{S'}$): $S'$ is finite free and $A$-Gorenstein. It also holds whenever $A$ is a DVR
  (\emph{no} hypothesis on $S'$). For $A$ normal, the identity therefore holds after localisation at every height-one prime.
\item In general $\cc=\mathrm{im}(\fb\to S')\cong\fb\cong S'^\vee$. So, if $S'$ is $A$-free (e.g.\ $S/\fb$ free, Lemma~19.5(c)), $\cc$ is a principal ideal of $S'$ if and only if
  (Gor$_{S'}$) holds.
\end{enumerate}
With $S=R$, $A=I_F$: under (Gor$_R$) and (Gor$_{R'}$), $\eta_F=\cc_q(F)\cdot\eta^B_F$, i.e.\ $(h_F)=\cc_q(F)$. \stproved{}
(Lean, namespace \pathtt{Langlandscape.CommAlg.ThmE}: (i) \lean{Langlandscape.CommAlg.ThmE.theta\_mem\_b}, \pathtt{mk\_theta\_mem\_ann}, \pathtt{eta\_eq\_span\_lambda'}, \lean{Langlandscape.CommAlg.ThmE.c\_eq\_orderIdeal};
(ii) \lean{Langlandscape.CommAlg.ThmE.eta\_eq\_c\_mul\_of\_unimodular};
(iii) Gorenstein half \lean{Langlandscape.CommAlg.ThmE.unimodular\_of\_isGorenstein}, \lean{Langlandscape.CommAlg.ThmE.eta\_eq\_c\_mul\_of\_isGorenstein}, DVR half \lean{Langlandscape.CommAlg.ThmE.eta\_eq\_c\_mul\_of\_dvr} (\pathtt{CommAlg/DVR.lean}, with \pathtt{exists\_unimodular\_of\_dvr});
(iv) \lean{Langlandscape.CommAlg.ThmE.c\_principal\_iff\_isGorenstein}.)
\stformalized{} for (i), (ii), (iii) (both halves) and (iv) for $S'$ $A$-free, under the hypotheses of Lemma~19.5 as the structure \texttt{NewOldData} (same qualifiers as there). \emph{Not formalized:} the localisation clause of (iii) for $A$ normal (height-one primes), the isomorphism $\mathrm{im}(\fb\to S')\cong\fb\cong S'^\vee$ as a standalone statement for general $S'$, and the final specialisation $S=R$, $A=I_F$ (the abstract statement is formalized, its instantiation to the Hecke rings is not).
\end{quoted}
\begin{outline}
Let $\ell$ be a trace of $S$, so $\ell(\theta s)=\lambda(s)$ (Lemma~19.3).
\begin{enumerate}[label=(\arabic*)]
\item (i) $\fa\subseteq\ker\lambda$ gives $\theta\fa=0$, so $\theta\in\Ann(\fa)=\fb$ (Lemma~19.5(a)); $\bar sx=\lambda'(\bar s)x$, so $x\in S'[\lambda']$; and
  $\eta(S,\lambda)=(\lambda(\theta))=(\lambda'(x))$. The key identity: for $\beta\in\fb$, $\lambda(\beta)=\ell(\theta\beta)=\Phi_\fb(\beta)(x)$, and
  $\Phi_\fb:\fb\xrightarrow{\sim}S'^\vee$ (Lemma~19.5(b)) is onto, so $\lambda(\fb)=\{\psi(x):\psi\in S'^\vee\}$. The seed's normalisation
  $\ell(\theta)=1$ is not needed. \stproved, \stformal
\item (ii) Write $x=h\theta'$, $h\in A$. Unimodularity means $\{\psi(\theta')\}=A$, so $\cc_q=hA$ and
  $\eta(S,\lambda)=(h\lambda'(\theta'))=\cc_q\cdot\eta(S',\lambda')$. \stproved, \stformal
\item (iii) Under (Gor$_{S'}$) with trace $\ell'$, Lemma~19.3 gives $\theta'$ with $\ell'(\theta')=\lambda'(1)=1$: unimodular. Over a DVR:
  $S'[\lambda']$ is saturated in the free module $S'$, hence a summand; it has rank one because $S_K=S'_K\times S''_K$ and
  $S_K[\lambda_K]=K\theta$ is one-dimensional by Lemma~19.3 for the $K$-Gorenstein $S_K$ (this uses (Gor$_S$)). For $A$ normal, localise
  at height-one primes (all three ideals localise; $\Ann$ of a f.g.\ ideal commutes with flat base change). \stproved; \stformal{} except
  the normal-$A$ clause.
\item (iv) $\fb\to S'$ is injective ($\fa\cap\fb=0$), so $\cc\cong\fb\cong S'^\vee$; principal iff $S'^\vee$ cyclic iff (for $S'$ finite free)
  $S'$ Gorenstein (Lemma~19.2). \stproved
\item Specialisation. $(R,\fa,\fb)$ over $I_F$ satisfies the hypotheses of Lemma~19.5 with $R''$ free (Lemma~20.4(c) below, \stpmc);
  (iii) gives $\eta_F=\cc_q(F)\eta^B_F$. Division needs $\eta^B_F\ne0$, which holds by multiplicity one for the primitive branch
  ($\T_\m\otimes\mathrm{Frac}\,I=\mathrm{Frac}\,I\oplus B$) \stcited{} \cite[(3.1), p.~18]{Hs} $=$ \cite[Cor.~3.7]{Hi88b}. Not formalized.
\end{enumerate}
\emph{Where the Gorenstein hypotheses enter} (main.pdf, Remark~19.9, p.~164). (Gor$_R$) is indispensable: it gives $\theta$, the
principality of $\eta_F$ and $\fb\cong R'^\vee$; without it nothing relates $\eta_F$ to $\fb$. (Gor$_{R'}$) is used only for unimodularity
of $\theta'$ and can be weakened to the condition of (ii); it is not needed at all over a DVR, i.e.\ at fixed weight. Sources:
(Gor$_R$) is \cite[\S3.3, p.~17--18]{Hs} \stcited{} under (CR), via \cite[Cor.~2, p.~482]{Wi95} and \cite[\S9 Prop.~2]{MW86}; (Gor$_{R'}$) is
\cite[Thm~4.5 and proof, p.~34]{Hs} \stcited{} for the definite algebra, transported by the JL isomorphism $\T'\cong\T^{\ord}_B$
(\cite[Rem.~7.8, p.~71]{Hs}). As a consistency check, Theorem~19.8 gives $\eta_F\subseteq\eta^B_F$, which is the inclusion of
\cite[Rem.~7.8]{Hs}. The toy algebra $S=\{(x,y)\in A^2:x\equiv y\bmod\varpi^t\}$ (main.pdf, Example~19.10, p.~165; Lean test
\pathtt{LanglandscapeTests/S06/Toy.lean}) realises every hypothesis with $\eta(S,\lambda)=\cc_q=(\varpi^t)$, $S'=A$.
\end{outline}

\paragraph{B. The old/new decomposition and the Tamagawa ideal (main.pdf, \S20).} These lemmas verify the hypotheses of
Theorem~19.8 for the Hecke rings and identify the specialisation of $\tq_q(F)$.

\begin{quoted}{Lemma}{20.4}{166}{lem:s06-decomp}{old/new decomposition}{chapters/s06/s06-m2.tex}\label{rpt:lem:s06-decomp}
Assume (H1). (a) $\T\otimes_\Lam K=(\T'\otimes K)\times(\T''\otimes K)$ with $K=\mathrm{Frac}\,\Lam$, $\fa\cap\fb=0$, and
$\T/\fa$, $\T/\fb$ are $\Lam$-torsion-free. (b) Let $\T(M)_\m$ be Hida's ordinary algebra of tame level $M$ localised at the
pull-back of $\m$, and $S_q\in\Lam^\times$ the element with $P_k(S_q)=q^{k-2}$ (trivial character). If $\bar\rho_\m$ is ramified at
$q$, then $\T''=0$. If $\bar\rho_\m$ is unramified at $q$ and $\T''\ne0$ (so that the pull-back $\m_M$ of $\m$ to $\T(M)$ is defined; if $\T''=0$ put $\T(M)_\m:=0$), $X^2-T_qX+qS_q$ has two roots in $\T(M)_\m$ that are distinct modulo the maximal ideal: they are
$\equiv U_q$ and $\equiv qU_q$ mod $\m$, and $q\not\equiv1$ by (H0). Let $\alpha_q$ be the root congruent to $U_q$. Then
$\T''\cong\T(M)_\m$ via $U_q\mapsto\alpha_q$. In particular $\T''$ is $\Lam$-free, and $T_q=U_q+qS_qU_q^{-1}$ lies in the image
of $\T$. The same holds at fixed weight $k$: $\T''_k\cong\T_k(M)_\m$, $U_q\mapsto\alpha_q$.
\revnote{F2: the old quotient is the localisation $\T(M)_\m$, not $\T(M)_\m[X]/(X^2-T_qX+qS_q)$ (which has twice the rank);
before: ``$\T''\cong\T(M)_\m[X]/(X^2-T_qX+qS_q)$, $X\mapsto U_q$; in particular $\T''$ is $\Lam$-free''.} (c) The same statements
hold after $\otimes_\Lam I_F$ for $R,R',R''$; in particular $(R,\fa,\fb)$ satisfies the hypotheses of Lemma~19.5 over
$A=I_F$, and $R''=R/\fb$ is $I_F$-free. \stpmc{} (proof below; cited inputs: Hida freeness and control, local--global
compatibility at $q$.)
\end{quoted}
\begin{outline}
Cited inputs: (I1) $\Lam$-freeness of $\T$, $\T(M)$ and control, \cite[Thm~3.1, Cor.~3.2, p.~250]{Hi86b} \stcited; (I2) local--global
compatibility at $\ell\ne p$ (unramified if $\ell\nmid N_f$; Steinberg shape with $a_\ell(f)^2=\ell^{k-2}$ if $\ell\,\|\,N_f$),
\cite[Thm~3.1]{DDT97}, \cite[Th.~(A)]{Car86} \stcited.
\begin{enumerate}[label=(\arabic*)]
\item (a) $\T_K$ is a finite product of local Artinian $\T_{K,\mathfrak n}$; each $\mathfrak n$ is an ordinary $\Lam$-adic eigensystem $\phi$ with a
  Galois representation $\rho_\phi$ (\cite[Thm~2.1]{Hi86a}, \stcited). Call $\mathfrak n$ $q$-new if $\rho_\phi$ is ramified at $q$, $q$-old
  otherwise; at a dense set of arithmetic points (I2) shows $q$-old $\iff$ from level $M$. Hence $\T_K=\T'_K\times\T''_K$ (``not both'' is the
  local dichotomy (I2), replacing a strong-multiplicity-one argument). $\fa\cap\fb$ kills a full-rank submodule of the free $\Lam$-module of
  cusp forms, hence is $0$ by faithfulness; torsion-freeness from $\T'\subseteq\End_\Lam(\mathbf S^{\rm new})$. \stpmc
\item (b) If $\T''\ne0$, a $q$-old $\phi$ at $\m$ forces $\bar\rho_\m$ unramified at $q$ (Brauer--Nesbitt, Chebotarev, (H2)). The old space is the image
  of $\mathbf S(M)^2$ under the two degeneracy maps, on which $U_q$ acts by $\big(\begin{smallmatrix}T_q&1\\-qS_q&0\end{smallmatrix}\big)$. Mod $\m$ the
  roots of $X^2-T_qX+qS_q$ are $u=\bar U_q$ and $qu$ (using $U_q^2=S_q$ on $q$-new forms, i.e.\ the existence of $F$ in the $\m$-component);
  they are distinct by $q\not\equiv1$ (H0). Hensel (discriminant a unit, $p\ne2$) splits the polynomial over $\T(M)_\m$; localising at $\m$
  keeps the factor $U_q\equiv u$; the resulting algebra contains $T_q=\alpha_q+qS_q\alpha_q^{-1}$, hence all of $\T(M)_\m$. The same argument
  at fixed weight. \stpmc
\item (c) $I_F$ is $\Lam$-free (H1), so $\otimes_\Lam I_F$ is exact and preserves every hypothesis of Lemma~19.5. \stproved{} from (a), (b).
\end{enumerate}
Hida's old ideal \cite[(3.4a--d)]{Hi86b} (old at any divisor of the tame level) is in general smaller than $\fb$ ($q$-old only); this is why
$\T''\cong\T(M)_\m$ is proved here rather than cited.
\end{outline}

\begin{quoted}{Lemma}{20.7}{169}{lem:s06-tam}{}{chapters/s06/s06-m2.tex}\label{rpt:lem:s06-tam}
(a) $\rho_F|_{I_q}$ is unipotent and factors through the $\Zp(1)$-quotient of tame inertia. So $\tq_q(F)$ is generated by the
entries of $\rho_F(\tau)-1$ for one generator $\tau$. (b) At an arithmetic point $P$ with $q$-new $f=F_P$,
$P(\tq_q(F))\,\widetilde\cO_P=(\varpi_P^{t_f(q)})$, where $\widetilde\cO_P$ is the normalisation of $I_F/P$ in its fraction field and $P(\tq_q(F))\widetilde\cO_P$ is the ideal generated by the image of $\tq_q(F)$ under $I_F\to I_F/P\subseteq\widetilde\cO_P$ (see the remark after the proof), and $t_f(q)$ is the Tamagawa exponent of
\cite[Def.~2.12]{KO} \cite[Def.~2.12,
p.~6]{KO} (= \cite[Def.~3.3]{PW} \cite[Def.~3.3, p.~9]{PW}),
computed with coefficient ring $\widetilde\cO_P$. (c) $t_f(q)<\infty$ because $f$ is $q$-new. \stpmc{} (proof below; inputs (I2) and
uniqueness of the lattice under (H2).)
\end{quoted}
\begin{outline}
(1) Unipotence: at a Zariski-dense set of arithmetic points, $\rho_f(\sigma)$ is unipotent by (I2); the intersection of these height-one
primes is $0$, so $(\rho_F(\sigma)-1)^2=0$. Wild inertia acts trivially ($q\in I_F^\times$); tame inertia is procyclic and its image
$\{1+aN:a\in\Zp\}$ is pro-$p$. (2) At $P$, $\mathcal T_F\otimes\widetilde\cO_P$ is the unique stable lattice up to scaling ((H2)), i.e.\
\cite{KO}'s $T_f$; $t_f(q)$ is the largest $t$ with $\rho_f(\tau)\equiv1\bmod\varpi_P^t$, so the ideal of entries is $(\varpi_P^{t_f(q)})$.
(3) $f$ $q$-new gives $\rho_f(\tau)\ne1$. \stpmc{} (inputs (I2) \stcited). \emph{The ring $\widetilde\cO_P$}: $I_F/P$ need not be normal, so one
passes to its normalisation, a DVR finite over $\cO$, with $I_F\to\widetilde\cO_P$ a \emph{local} homomorphism that is in general not surjective
(complete local rings are Nagata). This locality is what the squeeze in Theorem~21.5 uses. (A ramification index $e$ rescales both sides
of every identity equally.)
\end{outline}

\paragraph{C. The fixed-weight input and the $\Lam$-adic defect formula (main.pdf, \S21).}

\begin{quoted}{Theorem}{21.1}{170}{thm:s06-KO}{Kim--Ota}{chapters/s06/s06-m3.tex}\label{rpt:thm:s06-KO}
Let $p\ge5$ and $f\in S_k(\Gamma_0(N))$ a newform with $N=N^+N^-$, and assume: (1) $(N,p)=1$; (2) $\bar\rho_f|_{G_{\Q(\sqrt{p^*})}}$
is absolutely irreducible; (3) $2\le k\le p-1$; (4) $(N^+,N^-)=1$; (5) $N^-$ squarefree; (6) if $q\mid N^-$ and $q\equiv\pm1\bmod p$,
then $\bar\rho_f$ is ramified at $q$. Then $\ord_\lambda\eta_f(N)=\ord_\lambda\eta_f(N^+,N^-)+\sum_{q\mid N^-}t_f(q)$, where the
congruence ideals are taken in the full and in the $N^-$-new Hecke algebras of $S_k(\Gamma_0(N))$, localised at $\m_f$.
\stcited{} \cite[Thm~1.1]{KO}
\cite[Thm 1.1, p.~1;
Def.~2.11 ($N^-$-new congruence ideal), p.~6]{KO}. At $k=2$ it is \cite[Thm~6.8]{PW}
\cite[Thm 6.8, p.~21]{PW} with the definite self-pairing
$\xi_f$, under PW's (CR), whose (2) is the same $\pm1$ condition, and which also needs (CR)(1) ($\bar\rho_f$ surjective) and an odd number of primes in $N^-$ (harmless here: $N^-=q$) \cite[(CR), p.~2]{PW}.
\end{quoted}
\noindent \cite{KO}'s Hecke algebra is the full algebra on $H^1_{\text{ét},c}(Y_1(N),\mathcal F_k)[\langle\cdot\rangle-1]$ localised at $\m_f$,
which equals the image of the anemic algebra after localisation \cite[\S2.1.3, p.~5]{KO}; so it is our $\T_k(N)_\m$ (with $U_q$).

\begin{quoted}{Corollary}{21.2}{170}{cor:s06-fixed}{fixed-weight level lowering mod $\varpi^t$; replaces seed step (4)}{chapters/s06/s06-m3.tex}\label{rpt:cor:s06-fixed}
Let $f$ be a newform of level $N=qM$ and weight $2\le k\le p-1$, with $p\nmid N$, (H0), (H0$^-$), (H2) for $\bar\rho_f$, and $\m_f$ non-Eisenstein. Then $\cc_q(f)=(\varpi^{t_f(q)})$, and in particular $\cc_q(f)\subseteq(\varpi^{t_f(q)})$. \stpmc{}
\end{quoted}
\begin{outline}
\begin{enumerate}[label=(\arabic*)]
\item \emph{Gorenstein input.} $\T_k(N)_{\m_f}$ is $\cO$-Gorenstein: in general \cite[Prop.~5.1(1)]{KO}, citing \cite[Prop.~5.5]{Dim09},
  \stcited. In the application $f=F_P$ it is \emph{derived}: $R\otimes_{I_F,P}\widetilde\cO_P=\T_k(N)_\m\otimes\widetilde\cO_P$ (Lemma~21.4) is
  Gorenstein by Lemma~19.3 applied to (Gor$_R$); the argument is then run over $\widetilde\cO_P$ (flat over $\cO$). \cite[Cor.~5.4]{KO} is
  not used beyond $q\not\equiv1$ (it conflicts with \cite{BKM} at trivial primes).
\item \emph{Commutative algebra.} Lemma~20.4 at weight $k$ gives the hypotheses of Lemma~19.5 over the DVR $\cO$; Theorem~19.8(iii)
  over a DVR (no condition on the new quotient) gives $\eta_f(N)=\cc_q(f)\,\eta_f(M,q)$, where $\eta_f(M,q)=\eta(\T'_k,\lambda'_f)$ is literally
  \cite{KO}'s $N^-$-new congruence ideal \cite[Def.~2.11]{KO} for $N^-=q$, $N^+=M$. \stproved{} (given (1) and Lemma~20.4).
\item \emph{Comparison.} The hypotheses of Theorem~21.1 hold for $N^-=q$: (6) is exactly (H0) ($q\not\equiv1$) plus (H0$^-$). So
  $\ord_\varpi\eta_f(N)=\ord_\varpi\eta_f(M,q)+t_f(q)$; both ideals are non-zero ($f$ a newform), hence $\ord_\varpi\cc_q(f)=t_f(q)$.
  \stpmc{} (\cite[Thm~1.1]{KO}).
\end{enumerate}
So under (H0$^-$) the seed's step (4) needs no new $R=\T$ argument: it \emph{is} \cite[Thm~1.1]{KO} read through the fixed-weight
Theorem~19.8. Without (H0$^-$) the seed's direct $R=\T$ argument gives the inclusion (Prop.~23.2).
\end{outline}

\begin{quoted}{Lemma}{21.4}{171}{lem:s06-control}{control as an equality; step (3)}{chapters/s06/s06-m3.tex}\label{rpt:lem:s06-control}
Assume (H0), (H1), (H2), (Gor$_\T$) ($\T$ is $\Lam$-Gorenstein, see step (1) below), and let $P$ be an arithmetic prime of $I_F$
over $P_k$, $k\ge2$, with trivial character, $q$-new specialisation $f=F_P$ ($F_P$ is the $p$-stabilisation of a newform of level $N$, $F$ being primitive of tame level $N$), and (if $k=2$) the condition of (P2). Then
\begin{enumerate}[label=(\alph*)]
\item $\T/P_k\T\cong\T_k(N)_\m$, $\T''/P_k\T''\cong\T''_k\cong\T_k(M)_\m$, $\T'/P_k\T'\cong\T'_k$ (the $q$-new quotient at weight $k$), and
  $\fb\otimes_\Lam\Lam/P_k=\fb_k$, compatibly with the maps $\T\to\T',\T''$;
\item $P(\cc_q(F))\widetilde\cO_P=\cc_q(f)$ (equality, not only ``$\subseteq$'').
\end{enumerate}
\stpmc{} (cited: Hida control.)
\end{quoted}
\begin{outline}
Hida control \cite[Cor.~3.2, (3.3b), p.~250]{Hi86b} \stcited{} at levels $Np$ and $Mp$, followed by ordinary $p$-stabilisation
(main.pdf, Lemma~21.3, p.~171, \stpmc: for $k\ge3$ the ordinary part at level $Np$ is $p$-old since $a_p^2=p^{k-2}$ on $p$-new forms; for
$k=2$ this is the extra condition of (P2)), identifies $\T/P_k\T$ and $\T(M)_\m/P_k$. $\fb$: the sequence $0\to\fb\to\T\to\T''\to0$ splits
($\T''$ free), so it stays exact mod $P_k$. $\T'$: free by Lemma~19.5(c) (this needs (Gor$_\T$), which the seed did not name); the image of
$\fa$ kills $\fb_k$, so lies in $\fa_k$ by Lemma~19.5(a) at weight $k$, giving a surjection $\T'/P_k\T'\to\T'_k$ of free $\cO$-modules of equal
rank, hence an isomorphism. (b) follows since $\lambda_f$ factors through $\T/P_k\T$. \stpmc. The seed stated only ``$\subseteq$''; both
directions are used in Corollary~21.6.
\end{outline}

\begin{quoted}{Theorem}{21.5}{172}{thm:s06-defect}{$\Lam$-adic defect formula; seed Thm~defect, corrected}{chapters/s06/s06-m3.tex}\label{rpt:thm:s06-defect}
Assume (H0), (H1), (H2), (P) (with (P2) in its corrected form), and (H0$^-$) or (Lev) (Hyp.~20.1). Then
\[(h_F)=\cc_q(F)=\tq_q(F)\qquad\text{as principal ideals of }I_F.\]
\stpmc{} under (H0$^-$). Under (Lev) (this is where $q\equiv-1\bmod p$ with $\bar\rho$ unramified at $q$ is covered): \stpmc{} for $M=1$
and for $M>1$ in the cases (F1), (F1$'$) of Prop.~23.2 (for (F1$'$) in our sign-fixed reading of \cite[Thm~3.14]{KO}, not
\cite{KO}'s text). For $q\equiv-1\bmod p$ with $\bar\rho$ unramified at $q$ and outside (Lev) ($N(\bar\rho)>1$, or $\ell^2\mid M$ at a prime $\ell$ where $\bar\rho$ is unramified) the theorem is \stopen{} (Remark~25.1(i), Remark~23.8).

 Cited inputs: Hida freeness
and control \cite[Thm~3.1, Cor.~3.2]{Hi86b} \cite[Thm 3.1, Cor.~3.2
(3.3b), p.~250]{Hi86b}; (Gor$_R$) and (Gor$_{R'}$), with freeness of $H$ (\cite{Hs} as above, resting on \cite{Wi95}, \cite{MW86}, \cite[Cor.~8.11]{Hel}
\cite[Cor.~8.11, p.~31]{Hel04}); the $\Lam$-adic JL
(\pathtt{if:s01-targets}); and Thm~21.1 at the point $P$ of (P) under (H0$^-$), or Prop.~23.2 at $P$ under
(Lev) (resting on \cite[Thm~2.9, Lemma~2.4]{KO} for $M=1$, \cite[Thm~3.14]{KO} for (F1), (F1$'$)). (Im) is not used.
\end{quoted}
\begin{outline}
Let $P$ be the point of (P), $k=k_P$, $f=F_P$, $t=t_f(q)$, and $P:I_F\to\widetilde\cO_P$ the local homomorphism to the normalisation of
$I_F/P$.
\begin{enumerate}[label=(\arabic*)]
\item \emph{Gorenstein, and the combinatorial identity.} (Gor$_\T$) is \cite[\S3.3, p.~17--18]{Hs} \stcited{} under (CR), which holds by (H2)
  and (H0); (Gor$_R$) by base change (Lemma~19.3). (Gor$_{R'}$): \cite[Thm~4.5 and proof, p.~34]{Hs} \stcited{} for the definite
  $\T^{\ord}_{B,\m}$ under (CR,$\{q\}$), transported by the $\Lam$-adic JL isomorphism $\T'\cong\T^{\ord}_B$ (stage 01, \pathtt{if:s01-targets};
  \cite[Rem.~7.8]{Hs}), with the two $\Lam$-structures matched by $[u]\mapsto[u]^{-1}$ (classical versus adelic diamonds). Then
  Theorem~19.8(iii) and Lemma~20.4(c) give $\eta_F=\cc_q(F)\,\eta^B_F$, so $(h_F)=\cc_q(F)=(h)$ is principal. \stpmc
\item \emph{Galois half.} $\tq_q(F)\subseteq\cc_q(F)=(h)$ (Prop.~22.1). \stpmc
\item \emph{Control.} $P(\cc_q(F))\widetilde\cO_P=\cc_q(f)$ (Lemma~21.4). \stpmc{} (Hida control \stcited)
\item \emph{Fixed weight.} Under (H0$^-$): $\cc_q(f)=(\varpi_P^t)$ (Cor.~21.2, \stpmc{} via \cite[Thm~1.1]{KO}). Under (Lev): only
  $\cc_q(f)\subseteq(\varpi_P^t)$ (Prop.~23.2, cases (F0), (F1), (F1$'$), with the statuses there). The squeeze needs only the inclusion.
\item \emph{Squeeze, with $P$ local.} By Lemma~20.7(b),(c),
  \[(\varpi_P^t)=P(\tq_q(F))\widetilde\cO_P\subseteq P(\cc_q(F))\widetilde\cO_P=\cc_q(f)\subseteq(\varpi_P^t),\]
  so all are equal and $P(h)\ne0$. Write $\tq_q(F)=h\mathfrak d$ ($I_F$ a domain). Cancelling $P(h)$ in $\widetilde\cO_P$ gives
  $P(\mathfrak d)\widetilde\cO_P=\widetilde\cO_P$, so some $d\in\mathfrak d$ has $P(d)$ a unit; since $P$ is a \emph{local} homomorphism, $d\in I_F^\times$.
  Hence $\mathfrak d=I_F$ and $\tq_q(F)=(h)=\cc_q(F)=(h_F)$. \stproved{} as commutative algebra, given (1)--(4) and Lemma~20.7(b),(c) (\stpmc).
\end{enumerate}
Locality cannot be dropped: for $\Z_{(p)}\to\Q$, $\tq=(p)$, $h=1$ the squeeze fails. Surjectivity of $P$ is not needed (in Lean the
hypothesis is \texttt{[IsLocalHom P]} of \pathtt{CommAlg.squeeze}). \emph{Lean:} \pathtt{Defect.defect\_formula} is \stformal{} \emph{modulo}
the registered hypotheses \pathtt{Hyp.OldDecomp}, \pathtt{Hyp.HsiehGorFull}, \pathtt{Hyp.HsiehGor}, \pathtt{Hyp.LambdaJL},
\texttt{Hyp.TamSpec}, \pathtt{Hyp.HidaControlGL2}, \texttt{Hyp.KimOta} or \pathtt{Hyp.KimOtaRT}, and the Galois half
(\pathtt{Hyp.FreeRankTwo}, below); it is not an unconditional theorem. It is a statement about an abstract \pathtt{Defect.Setting};
Hida's Hecke algebra and the Galois representation are not constructed in Lean, and (H1), (H2), (P), (Lev) are not Lean arguments.
The toy instances of these hypotheses in the Lean tests are non-vacuity checks only, not evidence.
\end{outline}

\begin{quoted}{Corollary}{21.6}{173}{cor:s06-allweights}{all weights of the branch; seed Cor.~allweights}{chapters/s06/s06-m3.tex}\label{rpt:cor:s06-allweights}
Under the hypotheses of Thm~21.5, let $P'$ be any arithmetic point of weight $k'\ge3$ with trivial character
(including $k'\ge p$), with $q$-new $f'=F_{P'}$ ($F_{P'}$ is the $p$-stabilisation of a newform of level $N$, $F$ being primitive of tame level $N$). Then $\cc_q(f')=(\varpi_{P'}^{t_{f'}(q)})$ (in $\widetilde\cO_{P'}$) and
$\ord_{\varpi_{P'}}\eta_{f'}(N)-\ord_{\varpi_{P'}}\eta_{f'}(M,q)=t_{f'}(q)$, for the ring-theoretic congruence ideals of $\T_{k'}(N)_\m$ and of its
$q$-new quotient. \stpmc{} (with the status of the case of Thm~21.5 that is used; under (Lev), case (F1$'$), in our reading of \cite[Thm~3.14]{KO}).
\end{quoted}
\begin{outline}
From $\cc_q(F)=\tq_q(F)=(h)$ (Theorem~21.5): Lemma~21.4(b) at $P'$ (valid for every $k'\ge3$, including $k'\ge p$) and Lemma~20.7(b) at $P'$
give $\cc_q(f')=P'(\cc_q(F))\widetilde\cO_{P'}=P'(\tq_q(F))\widetilde\cO_{P'}=(\varpi_{P'}^{t_{f'}(q)})$. Then $\T_{k'}(N)_\m\otimes\widetilde\cO_{P'}$ is Gorenstein by
base change of (Gor$_R$), and Theorem~19.8(iii) over the DVR $\widetilde\cO_{P'}$ gives $\eta_{f'}(N)=\cc_q(f')\eta_{f'}(M,q)$; both are non-zero
($f'$ a newform); take $\ord$. \stpmc. No $R=\T$ theorem is used at $P'$, only Hida control, $p$-stabilisation and Lemma~20.7. For
$k'\ge p$ this goes beyond \cite{KO} (which has $k\le p-1$); this, together with $(h_F)=\cc_q(F)$ on the new branch, is the new content.
Lean: \pathtt{Defect.cor\_allweights} gives the ideal identity; the $\ord$-formula additionally needs $\eta_{f'}(N)\ne0\ne\eta_{f'}(M,q)$, which
is not assumed in Lean.
\end{outline}

\paragraph{How the integral identity relates to the literature, as far as the chapter goes.} At a point $P$ the stage proves
$\eta_f(N)=\cc_q(f)\,\eta_f(M,q)$ from commutative algebra alone (Theorem~19.8(iii) over a DVR, given Gorensteinness of $\T_k(N)_\m$), and
$\eta_f(M,q)$ is the congruence ideal of the $q$-new quotient, which is the definite side through the JL isomorphism $\T'\cong\T^{\ord}_B$
(stages 01, 04). The arithmetic content is $\cc_q(f)=(\varpi^{t_f(q)})$, and the chapter relates it to the literature only as follows.
(i) Under (H0$^-$) it is equivalent, via the identity above, to the Kim--Ota formula \cite[Thm~1.1]{KO} for $N^-=q$ (Corollary~21.2); at
$k=2$ that formula is Pollack--Weston \cite[Thm~6.8]{PW}. (ii) The $R=\T$ theorems (\cite[Thm~0.2]{DFG}, \cite[Thm~6.6]{Dim09},
\cite[Thm~2.9]{KO}) and B\"ockle's complete-intersection presentation \cite[Thm~3.12]{KO} (citing \cite{Boc07}) enter only inside the proof of
\cite{KO}, or (outside (H0$^-$)) through \cite[Thm~3.14]{KO} in Prop.~23.2. (iii) Diamond \cite[Thm~6B, p.~80]{Dia89} computes the same
$\Lam$-adic old/new congruence module along the \emph{old} branch; the stage reads it along the \emph{new} branch through $\lambda'$.
(iv) The inclusion $\eta_F\subseteq\eta^B_F$ given by Theorem~19.8 is the inclusion of \cite[Rem.~7.8]{Hs}. (v) At trivial primes
(excluded by (H0)), \cite[Thm~1.2]{BKM} computes a Wiles defect at fixed weight; a comparison with $\cc$ is not attempted (main.pdf,
Remark~25.1(iii)). The chapter makes no further comparison; in particular our proofs do not use Wiles' numerical criterion or
Ribet--Takahashi; they appear only as heuristic context in \S\ref{ss:s00-oneway}.

\paragraph{D. The Galois half (main.pdf, \S22).} The box below contains, verbatim, the binding disclosure sentence on the Lean status
of the Galois half.

\begin{quoted}{Proposition}{22.1}{174}{prop:s06-galois}{Galois comparison; seed Prop.~galois}{chapters/s06/s06-m4.tex}\label{rpt:prop:s06-galois}
Assume (H0), (H1), (H2). Then $\tq_q(F)\subseteq\cc_q(F)$. \stpmc{}
(Lean: Given \pathtt{Hyp.ResidualBurnside}, \pathtt{Hyp.FreeRankTwo} is equivalent to the inclusion $\tq_q\subseteq\cc_q$; both directions are proved in Lean (\pathtt{Defect.galois}, \pathtt{Defect.freeRankTwo\_of\_le}, packaged as \pathtt{Defect.freeRankTwo\_iff\_le}). The Galois half is therefore not formalized independently: its arithmetic content (that the old-form representation $\rho_M$ exists at level $M$, steps 1--3 of Prop.~22.1) is the cited hypothesis. \pathtt{Hyp.ResidualBurnside} is discharged only given absolute irreducibility, (H2).)
\par\smallskip\noindent{\scriptsize\itshape [Status paragraph following the statement in the chapter:]}\par
\noindent\emph{Cited inputs.}
(1) $H$ is free of rank two over $\T$, and $H(M)$ is free of rank two over $\T(M)_\m$ (tame level $M$):
\cite[\S7.2, p.~68]{Hs} \cite[\S7.2, proof of Lemma 7.4, p.~68]{Hs}
\stcited, citing \cite[Cor.~1, 2, p.~482]{Wi95}
\cite[Thm 2.1, Cor.~1, Cor.~2, p.~482]{Wi95}, under (CR) (which (H2)
implies, Remark~20.2). The Galois action on $H$ (resp.\ $H(M)$) is unramified outside $Np$ (resp.\ $Mp$), and
$\tr(\Frob_\ell\mid H)=T_\ell$ in $\T$ for $\ell\nmid Np$ (Eichler--Shimura; for each classical specialisation this is
\cite[Thm 3.1(a), p.~85]{DDT97}) \stcited.
(2) $H(M)$ is unramified at $q$, since $q\nmid Mp$ ($X_1(Mp^r)$ has good reduction at $q$; for each classical specialisation this is
again \cite[Thm 3.1(a), p.~85]{DDT97}) \stknown.
(3) Chebotarev: for a Galois extension of $\Q$ unramified outside a finite set $S$, the Frobenius elements at primes $\notin S$ are
dense \cite[Thm 2.3, p.~52]{DDT97} \stknown.
\end{quoted}
\begin{outline}
If $\bar\rho_\m$ is ramified at $q$, then $\T''=0$, $\cc_q(F)=I_F$ and there is nothing to prove. Otherwise let $A:=I_F/\cc_q(F)\ne0$, a local ring.
\begin{enumerate}[label=(\arabic*)]
\item $\lambda\bmod\cc_q(F)$ kills $\fb R$, so factors through $R''\cong\T(M)_\m\otimes I_F$ (Lemma~20.4(b),(c)), giving $\nu:\T(M)_\m\to A$ with
  $\nu(T_\ell)=\lambda_F(T_\ell)\bmod\cc_q(F)$ for $\ell\nmid Np$. \stpmc{} (Lemma~20.4)
\item Two representations over $A$: $\rho_A=\rho_F\otimes A$ and $\rho_M=H(M)\otimes_{\nu}A$ (input (1)); their traces agree at
  $\Frob_\ell$, $\ell\nmid Np$. \stcited{} (input (1))
\item Chebotarev and continuity ($A$ Hausdorff): $\tr\rho_A=\tr\rho_M$ on $G_\Q$. \stknown{} (input (3))
\item For $\sigma\in I_q$, $\rho_M$ is unramified at $q$ (input (2)), so $\tr\rho_A(g\sigma)=\tr\rho_A(g)$ for all $g$.
\item Burnside in dimension two (main.pdf, Lemma~22.2, p.~174: for $\bar\rho$ absolutely irreducible and $A$ local, $A[G]\to M_2(A)$ is
  surjective; \stproved, \stformal{} as \pathtt{CommAlg.surjective\_groupAlgebra\_of\_absIrred}) and the perfect trace form (main.pdf,
  Lemma~22.3: $\tr\rho(g\sigma)=\tr\rho(g)$ for all $g$ forces $\rho(\sigma)=1$; \stproved, \stformal{} as
  \pathtt{CommAlg.eq\_one\_of\_trace\_mul\_eq}) give $\rho_A(\sigma)=1$, i.e.\ the entries of $\rho_F(\sigma)-1$ lie in $\cc_q(F)$.
\end{enumerate}
(H0$^-$) is not used: $q\equiv-1\bmod p$ is allowed here. (H0) is used only through $\T''\cong\T(M)_\m$. (Im) is not used. In Lean, steps
1--3 are the content of \pathtt{Hyp.FreeRankTwo}; the binding wording is: \pathtt{Defect.defect\_formula} is formalized modulo the
registered \texttt{Hyp.*}, including the Galois half (\pathtt{Hyp.FreeRankTwo}); it is not an unconditional theorem.
\end{outline}

\paragraph{E. The case $q\equiv-1\bmod p$ (main.pdf, \S23).} Corollary~21.2 needs (H0$^-$): \cite[Thm~1.1(6)]{KO} and \cite[(CR)(2)]{PW}
both exclude $q\equiv\pm1$ with $\bar\rho$ unramified at $q$ (main.pdf, Remark~23.5(a), \stcited). The precise place in \cite{KO} is the
L-value step \cite[Prop.~2.15(1), p.~8]{KO}: passing from the imprimitive to the primitive congruence ideal multiplies by
$L_q(\mathrm{ad}^0f,1)^{-1}=1-q^{-2}$, a $p$-unit iff $q\not\equiv\pm1$ (Remark~23.5(b), \stcited); this is a property of their method, not of
the statement $\cc_q(f)=(\varpi^{t_f(q)})$. The squeeze needs only the inclusion, and the seed's direct $R=\T$ argument gives it with no
condition at $q$ beyond (H0):

\begin{quoted}{Proposition}{23.2}{176}{prop:s06-levM}{the inclusion without (H0$^-$); the seed's step (4)}{chapters/s06/s06-m5.tex}\label{rpt:prop:s06-levM}
Let $f$ be as above, with (H0) and (Lev$'_M$) (Hyp.~23.1), and assume that $f$ is $p$-ordinary, i.e.\ $a_p(f)\in\cO^\times$
(automatic for $f=F_P$ in Thm~21.5). Then $\cc_q(f)\subseteq(\varpi^{t_f(q)})$ in each of the following cases, in which
(Lev$'_M$) holds automatically, $N(\bar\rho)=1$, $M_0=M$ and $\Sigma^+=P_{\bar\rho}=\emptyset$:
\begin{itemize}
\item[(F0)] $M=1$. \stpmc{} (Remark~23.10).
\item[(F1)] $M>1$ squarefree, $\bar\rho$ unramified at every $\ell\mid M$, and $\ell\not\equiv-1\bmod p$ for every $\ell\mid M$.
  \stpmc
\item[(F1$'$)] as (F1), but some $\ell\mid M$ has $\ell\equiv-1\bmod p$. \stpmc{} in the sign-fixed reading of \cite[Thm~3.14]{KO}
  (Remark~23.4), which is our reading, not \cite{KO}'s text (F2 part B and audit part 2: confirmed as a reading).

\end{itemize}
No condition on $q\bmod p$ beyond (H0) is used; in particular $q\equiv-1\bmod p$ is allowed. For $N(\bar\rho)>1$ (under (Lev$'_M$), with
or without (L3)) the inclusion is \stopen: the proof below does not apply when $\bar\rho$ is not twist-minimal (counterexample B2 of F2,
Remark~23.9), and the inclusion itself is neither proved nor disproved there.
\end{quoted}
\begin{outline}
Here $f$ is a $q$-new newform of level $N=qM$, weight $2\le k\le p-1$, trivial character, $p\nmid N$, (H2), $t=t_f(q)$. In (F0)--(F1$'$),
$N(\bar\rho)=1$, so every item concerning $\ell\mid N(\bar\rho)$ is void.
\begin{enumerate}[label=(\arabic*)]
\setcounter{enumi}{-1}
\item If $\bar\rho$ is ramified at $q$, then $t=0$ \cite[Def.~2.12]{KO} \stcited. Assume $\bar\rho$ unramified at $q$.
\item \emph{Reduction.} A ring map $\phi:\T''_k\to\cO/\varpi^t$ with $\phi\circ\pi''=\lambda_f\bmod\varpi^t$ gives $\lambda_f(\fb_k)\subseteq\varpi^t\cO$. \stproved
\item \emph{Old quotient.} $\T''_k\cong\T_k(M)_\m$, $U_q\mapsto\alpha_q$ (Lemma~20.4(b)); with $a_q(f)=\varepsilon c$, $c=q^{(k-2)/2}$, the two
  $U_q$-roots are $\equiv\varepsilon c,\varepsilon qc$, distinct by (H0) (at $q\equiv-1$ they are $\equiv\pm\varepsilon c$, still distinct). \stpmc
\item \emph{The point.} $\rho_t=\rho_f\bmod\varpi^t$ is a deformation of $\bar\rho$ satisfying the local conditions of \cite[Def.~3.9]{KO}:
  Fontaine--Laffaille at $p$ (\cite[\S3.2.1]{KO}, citing \cite{DFG}, \stcited); unramified at $q$ \emph{by the definition of $t$}; $\ell$-new
  (Steinberg of sign $\varepsilon_\ell$) at $\ell\in\Sigma^-$ by (L1) and local--global compatibility. So $\rho_t$ gives
  $x:R^{\Sigma^-}_{\Sigma^+}\to\cO/\varpi^t$. \stpmc
\item \emph{$R=\T$.} \cite[Thm~3.14, p.~12]{KO} with $N^-=M_0$, $N^+=N(\bar\rho)$ (level $rM$, $r$ an auxiliary prime) gives
  $\pi:R^{\Sigma^-}_{\Sigma^+}\xrightarrow{\sim}A=\T(rM)^{M_0\text{-new}}_{\m_{\Sigma^+}}$, induced by the Galois representation over $A$. It needs no
  condition $q\not\equiv\pm1$, and $q\notin\Sigma$: the deformation problem is unramified at $q$, so \cite[Prop.~2.15(1)]{KO} plays no role.
  \stcited{} for (F0), (F1); for (F1$'$) \stpmc{} in the sign-fixed reading below. Put
  $\psi:=x\circ\pi^{-1}\circ(\T_k(M)_\m\twoheadrightarrow B=A)$, so $\psi(\tau_\sigma)=\tr\rho_t(\sigma)$.
\item \emph{Matching.} By main.pdf, Lemma~23.3 (the full algebra $B=\T_k(M)^{M_0\text{-new}}_\m$ equals the anemic one $A$, with explicit
  trace formulas for $T_\ell$, $T_p$ and $U_\ell$; \stpmc{} when $N(\bar\rho)=1$, citing local--global compatibility, strong multiplicity
  one, Chebotarev, Deligne's theorem \cite[Thm~2.5]{Edi92} and Fontaine/Edixhoven \cite[Thm~2.6]{Edi92}, the last two not page-checked),
  $\psi(X')=\lambda_f(X)\bmod\varpi^t$ for each generator: $T_\ell$ by traces; $T_p$ through the unit root (all forms in $\m$ are ordinary);
  $U_\ell=a_\ell(f)$ for $\ell\in\Sigma^-$; and $U_q$: $\psi(T_q)=\varepsilon(q+1)c$, so $\psi(\alpha_q)$ is the root $\equiv\varepsilon c$, and by
  unit difference (step 2) $\psi(\alpha_q)=a_q(f)$, with no Hensel lemma over $\cO/\varpi^t$. So $\phi$ of step (1) exists. \stpmc
\end{enumerate}
\emph{The sign at $\ell\equiv-1\bmod p$} (main.pdf, Remark~23.4, p.~180, repair R4). For $\ell\in\Sigma^-$ the Steinberg lifts of the two
signs have residual Frobenius eigenvalues $\{\pm c_\ell,\pm\ell c_\ell\}$, which coincide iff $\ell\equiv-1\bmod p$. For $\ell\not\equiv-1$,
$\bar\rho|_{G_{\Q_\ell}}$ fixes the sign and \cite[Thm~3.14]{KO} applies as stated (\stcited). For $\ell\equiv-1$ both signs lift the same
$\bar\rho|_{G_{\Q_\ell}}$; \cite{KO}'s normal form fixes one sign while the anemic localisation contains both, and \cite{KO}'s proof does not
address this (BKM \cite[\S4.2, \S5]{BKM} fix the sign with $f$). The \emph{sign-fixed reading} takes the normal form of the sign of $f$ and
the factor of $A$ with $U_\ell=\varepsilon_\ell c_\ell$; \cite{KO}'s proof then goes through verbatim (finiteness \cite[Thm~2.9]{KO}, complete
intersection \cite[Lemma~3.8, Cor.~3.13]{KO}, flat and reduced \cite[Lemmas~3.15, 3.16]{KO}, existence and injectivity of $\pi$). \stpmc;
this is our reading, not \cite{KO}'s text.

\emph{Status} (main.pdf, Remark~23.8, p.~182). The team lifted the provisional marks on (F1) and (F1$'$) on 15 November 2026, on the basis
of two independent readings (an independent reviewer with a separate account, and the stage-06 chapter audit), both modulo
\cite[Thm~3.14]{KO}, local--global compatibility, Deligne's theorem and Fontaine/Edixhoven, and both confirming (F1$'$) as a reading,
not as \cite{KO}'s text. There is no numerical test of (F1) or (F1$'$): it needs $\bar\rho$ with $N(\bar\rho)=1$, i.e.\ from level-one forms
of weight $\ge12$ (e.g.\ $q=1741$ at $p=13$), which is out of reach here. A status \emph{proved} without cited input is not in reach:
the $R=\T$ theorem is cited.

\emph{The open case $N(\bar\rho)>1$} (main.pdf, Remark~23.9, p.~182). Counterexample B2 ($p=13$, $k=12$, $\bar\rho=\bar\rho_\Delta\otimes\bar\chi_{43}$,
$M=43^2$; the level-raising congruence at $43$ \stcomp{} from the $q$-expansion of $\Delta$): there is a newform $g=h\otimes\chi_{43}$ in
$\mathcal N$ that is not \cite{KO}-minimal, so \cite{KO}'s minimal ring has no point at $g$ while $A$ has its eigensystem; hence
$\pi:R\cong A$ of step (4) is false there ((L3) fails). B2 refutes \emph{proof step (4)}, not the inclusion $\cc_q(f)\subseteq(\varpi^t)$, and
not the matching Lemma~23.3 (B2 is in its case ($\delta$), which holds). Under (L3), step (4) rests on \cite[Lemma~2.4]{KO} (via
\cite{Dim09}, not read). A route under (L3) is \stsketch{} only; the case $N(\bar\rho)>1$ is \stopen. Numerical evidence in the open
region: the inclusion holds with equality in $5$ non-vacuous cases with $t=2$, all with $N(\bar\rho)>1$ \stcomp{}
[\res{s06-xcheck/H3.json}]; this is not a test of (F1).
\end{outline}

\paragraph{F. The Betti incarnation of $\cc_q$ (main.pdf, \S24).} A side result linking $\cc_q$ to Ribet's level-raising module of
stage 03, used for the numerical cross-checks.

\begin{quoted}{Proposition}{24.1}{184}{prop:s06-betti}{Betti incarnation of $\cc_q$}{chapters/s06/s06-m7.tex}\label{rpt:prop:s06-betti}
Let $l=p>k\ge4$ (the range of \pathtt{if:s03-amalgam}), $q\not\equiv-1\bmod p$, and let $\m$ be non-Eisenstein for $\T_k(M)$ acting on
$H^1(\Gamma_0(M),V_k)$, with $H^1(\Gamma_0(M),V_k)_\m$ free of rank two over $\T_k(M)_\m$. Then:
\begin{enumerate}[label=(\alph*)]
\item $\coker(\Psi^*\Psi)_\m\cong\big(\T_k(M)_\m/(\delta)\big)^{\oplus2}$, where $\delta:=\det(\Psi^*\Psi)=(q+1)^2-q^{2-k}T_q^2$.
  Since $p\nmid2(q+1)$, $(\delta)=(T_q-\epsilon(q+1)q^{(k-2)/2})$ for the unique sign $\epsilon$ with this element in $\m$.
\item $\T_k(N)_{\m}/(\fa+\fb)\cong\T_k(M)_\m/(T_q-\epsilon(q+1)q^{(k-2)/2})$ (Ribet's level-raising module, with Ihara's lemma).
\item If $\rk_\cO\T'_{k,\m}=1$, then $\ord_\varpi\cc_q(f)=\tfrac12\,\mathrm{length}_\cO\coker(\Psi^*\Psi)_\m$. In general only
  $\ord_\varpi\cc_q(f)\le\tfrac12\mathrm{length}$ holds.
\end{enumerate}
(a) \stproved{} (linear algebra, Lemma~24.2). (b) \stpmc{} (Ihara's lemma in weight $k$ and the cohomological
congruence module: \cite[Lemma~3.2 and \S2]{Dia91}, \cite[\S3]{Dia89}). (c) \stpmc{} (from (a), (b)).
Here: in (a), if neither sign gives an element of $\m$, then $\delta\in\T_k(M)_\m^\times$ and both sides are $0$ for either sign;
in (b), $\T_k(N)_\m$ is the localisation of the weight-$k$ Hecke algebra of $\Gamma_0(N)$ (with $U_q$; $N=qM$) at the maximal ideal $\m_N$
generated by $\m$ (on the $T_\ell$, $\ell\nmid Np$, and $U_\ell$, $\ell\mid M$) and $U_q-\epsilon c$, and $\fa,\fb$ are its $q$-new and
$q$-old kernels; in (c), $f$ is a $q$-new eigenform in $\m_N$, $\T'_{k,\m}$ is the $q$-new quotient of $\T_k(N)_\m$, $\cc_q(f)$ is as in
Def.~20.5, and (H0) is assumed. \revnote{S06-M7: the degenerate case ``no sign'' in (a) and the maximal ideal
$\m_N$ of (b) are made explicit, and (H0) ($p\nmid\varphi(N)$, $q\not\equiv1$) is used in (b); statement otherwise unchanged.}
\end{quoted}
\begin{outline}
From stage 03 only the shape $\Psi^*\Psi=\big(\begin{smallmatrix}q+1&-T'_q\\-T''_q&q+1\end{smallmatrix}\big)$, $T'_qT''_q=q^{2-k}T_q^2$, is used.
(a) For a unit $u$, $\coker\big(\begin{smallmatrix}u&-a\\-b&u\end{smallmatrix}\big)$ on $H^2$ is $H/(u^2-ab)H$ by elementary row/column
operations (main.pdf, Lemma~24.2, \stproved; \stformal{} for $H=T$ only, \pathtt{CommAlg.TwoByTwo.cokerEquivOfIsUnit}); $u=q+1$ is a unit by
$q\not\equiv-1$, and $\delta=-q^{2-k}(T_q-(q+1)c)(T_q+(q+1)c)$ with factors differing by a unit. \stproved.
(b) $\supseteq$: $U_q^2-q^{k-2}\in\fa$ by (I2). $\subseteq$: Ihara's lemma in weight $k$ \cite[Lemma~3.2, p.~210]{Dia91} \stcited{} and the
cohomological congruence module $\Omega\cong\ker(U_q^2-q^{k-2})$ \cite[\S2]{Dia91} \stcited; after non-Eisenstein localisation
$H^1=H^1_P$ (\stknown{} as standard, but no source located in the project literature: open point N5; (b) is \stpmc{} modulo this unsourced input as well); $T_q^2-(q+1)^2q^{k-2}$ is a non-zero-divisor by
Ramanujan--Petersson \cite[Thm~5.6]{Del69} \stcited; then $\Omega_{\m_N}\cong(T/(\alpha_q^2-q^{k-2}))^2$ and $\fa$ kills it. \stpmc.
(c) $\T'_k/\cc\cong\T_k/(\fa+\fb)$ surjects onto $\cO/\cc_q(f)$; equality when $\T'_k=\cO$. \stpmc. For $q\equiv-1$ (a) fails, and every use
for $k\ge p$ is \stexp{} (main.pdf, Remark~24.3).
\end{outline}

\paragraph{G. Several ramified primes (main.pdf, \S25).}

\begin{quoted}{Conjecture}{25.2}{187}{conj:s06-several}{several ramified primes; seed \S10, last paragraph}{chapters/s06/s06-m8.tex}\label{rpt:conj:s06-several}
Let $D=q_1\cdots q_r$ ($r$ odd), with each $q_i$ satisfying (H0), (H0$^-$) and $F$ $D$-new. Let
$\T=\T^{(0)}\twoheadrightarrow\T^{(1)}\twoheadrightarrow\cdots\twoheadrightarrow\T^{(r)}=\T^{D\text{-new}}\cong\T^{\ord}_{B_D}$ be the
$q_1\cdots q_i$-new quotients, and let $\cc^{(i)}_{q_i}(F)$ be the old/new congruence ideal of $\T^{(i-1)}\twoheadrightarrow\T^{(i)}$ at
$q_i$. Then $(h_F)=\prod_i\cc^{(i)}_{q_i}(F)=\prod_i\tq_{q_i}(F)$. \stconj{} If every $\T^{(i)}_\m$ is Gorenstein, the first
equality follows by iterating Thm~19.8 \stsketch. Whether the intermediate quotients are Gorenstein is \stopen{} (Hsieh
covers $i=0$ and $i=r$). The second equality would follow from Prop.~22.1 at each step and \cite[Thm~1.1]{KO} with
$N^-=D$, by the squeeze of step (5) \stsketch. Note that $q_i$ is $q_1\cdots q_{i-1}$-new-ramified in the intermediate algebras,
which needs the relative version of Lemma~20.4 at each step.
\end{quoted}
\noindent For $r=3$ (main.pdf, Remark~25.3, \stsketch): step $i$ is Theorem~19.8 for $S_{i-1}\twoheadrightarrow S_i$; the intermediate
algebras $S_1,S_2$ are each the ``$S$'' of one step, so their Gorenstein property is indispensable, and this is exactly what is missing
(for general odd $r$: the $r-1$ intermediate quotients). Given that, the product squeeze at one point with \cite[Thm~1.1]{KO} for $N^-=D$
yields all factorwise equalities $\cc^{(i)}=\tq_{q_i}$, since $I_F$ is local.

\noindent\emph{Secondary results not boxed here:} Lemma~21.3 ($p$-stabilisation, \stpmc), Lemma~22.2 (Burnside, \stproved, \stformal),
Lemma~22.3 (trace form, \stproved, \stformal), Lemma~23.3 ($B=A$, \stpmc{} for $N(\bar\rho)=1$), Lemma~24.2 ($2\times2$ cokernels, \stproved,
\stformal{} for $H=T$), Remark~23.5 (the case $q\equiv-1$, \stcited/\stpmc), Remark~25.1 (open cases, \stopen); all are used in the
outlines above with these statuses.

\subsubsection*{What this stage does not prove}
\begin{itemize}
\item $q\equiv-1\bmod p$ with $\bar\rho$ unramified at $q$, outside (Lev): $N(\bar\rho)>1$, and non-squarefree $M$ ($\ell^2\mid M$ at a prime
  where $\bar\rho$ is unramified), are \stopen{} (main.pdf, Remarks~23.9, 25.1(i)). Counterexample B2 refutes proof step (4) there, not the
  inclusion. (F1) and (F1$'$) are numerically untested; (F1$'$) rests on our sign-fixed reading R4 of \cite[Thm~3.14]{KO}.
\item Branches whose only weight-two points are $p$-new, or where the extra condition of (P2) fails: the Galois half would need \cite{KO} at
  level $Np$ \stopen. Trivial primes $q\equiv1\bmod p$: (Gor$_{R'}$) fails \cite[Thm~1.2]{BKM} \stcited, so $\cc$ is not principal; excluded by
  (H0). Dropping (Im) (probably vacuous) is not carried out (main.pdf, Remark~25.1).
\item Several ramified primes: Conjecture~25.2 is \stconj; the Gorenstein property of the intermediate new quotients is \stopen.
\item The fixed-weight and $\Lam$-adic Gorenstein properties, Hida freeness and control, freeness of $H$, the Galois representations, and
  the $R=\T$ theorems are cited, not proved; hence every arithmetic statement of the stage is at best \stpmc.
\item Lean: the commutative-algebra layer is formalized (qualifiers as in the boxes); not formalized are the normal-$A$ clause of
  Theorem~19.8(iii), the standalone $\mathrm{im}(\fb\to S')\cong\fb\cong S'^\vee$, the CRT step of Lemma~19.6, the base-change sentence of
  Lemma~19.5, and the specialisation of Theorem~19.8 to Hecke rings. The Defect layer is formalized only modulo the registered
  \texttt{Hyp.*}, including \pathtt{Hyp.FreeRankTwo}; eight of the nine stage-06 hypotheses are open (\pathtt{Hyp.ResidualBurnside} is
  conditionally discharged, given absolute irreducibility).
\item Prop.~24.1 for $k\ge p$ is \stexp; the input $H^1=H^1_P$ after non-Eisenstein localisation has no located source.
\end{itemize}
